% Compile from this directory with: pdflatex project_17.tex (twice).
\documentclass[11pt,a4paper]{article}
\usepackage[T1]{fontenc}
\usepackage[utf8]{inputenc}
\usepackage[margin=22mm]{geometry}
\usepackage{lmodern}
\usepackage{amsmath,amssymb}
\usepackage{enumitem}
\usepackage{xcolor}
\usepackage{microtype}
\usepackage{hyperref}
\definecolor{epflred}{HTML}{E3000B}
\hypersetup{colorlinks=true,urlcolor=epflred,linkcolor=epflred,
  pdftitle={ENG-654 Project 17: Global IK-Map Planning for KUKA iiwa 7}}
\setlength{\parindent}{0pt}
\setlength{\parskip}{6pt}
\setlist[itemize]{leftmargin=*,itemsep=5pt,topsep=4pt}
\urlstyle{tt}

\begin{document}
{\small\textbf{EPFL \textbar\ ENG-654}\hfill Kinematics-Grounded Motion Planning for Robots}
\vspace{5mm}

{\color{epflred}\Large\bfseries Project 17}

{\LARGE\bfseries Global IK-Map Planning for KUKA iiwa 7\par}

\textbf{Format:} Group of two students; simulation-based project.\\
\textbf{Course foundations:} Lectures 4 and 8; Exercise 4.

\section*{Motivation}
Choosing a redundant angle independently at every pose can produce
apparently feasible configurations that do not connect into a trajectory.
A global view of path progress and redundancy exposes which IK families
survive and whether a constant posture parameter must be replaced by a
varying schedule.

\section*{Task}
\textbf{Overall objectives.} Build classified IK feasibility charts for one iiwa 7 full-pose path,
then compare constant-$q_3$ choices with a continuous varying-$q_3$ schedule
through validated feasible connections.

\begin{itemize}
  \item \textbf{Validate the path and model.} Inspect the URDF, D--H conventions,
physical limits, and tool frame. Read the supplied pose CSV and check
quaternion normalization, component order, and pose continuity. Reuse the
course analytical IK and verify its outputs by full forward kinematics.

  \item \textbf{Enumerate and classify the chart.} Let $\lambda=q_3$ and evaluate
each pose sample over a sorted redundancy grid. Retain the supplied analytical
labels for up to eight fixed-angle candidates, including missing-solution
slots. Explain that labels in a fixed-angle slice are not by themselves
proof of global component membership.

  \item \textbf{Apply explicit feasibility checks.} Record position/orientation
residuals, joint-limit margins, full-Jacobian regularity, and reduced-chart
regularity. Produce separate path-sample versus redundant-angle panels for
each label with failure reasons. Plot the smallest singular value of the
scaled $6\times7$ Jacobian separately from the square reduced determinant.

  \item \textbf{Compare constant and varying schedules.} Test each sampled
constant-$\lambda$ row for complete-path continuation. Build a layered graph
or dynamic program over neighboring feasible candidates to minimize joint
travel. An edge is accepted only after intermediate poses are re-solved and
checked; proximity in the image is not sufficient.

  \item \textbf{Validate the selected trajectory.} Re-solve the final schedule on
a denser path grid and compare its completion, travel, and minimum margins
with the best constant-angle candidate. Preserve actual joint continuity
within physical limits. Report any topology or connection that remains
unresolved at the chosen resolution.
\end{itemize}

\newpage
\section*{Specifics and scope}
Use one course pose path and reuse the Exercise 4 enumeration if
available. Its reference grid has 175 negative and 175 positive redundancy
angles, excluding zero, with bounds documented from the assignment model.
For development, a smaller grid is acceptable; perform the final comparison
on the reference grid or justify an adaptive equivalent. Limit refinement
to one critical corridor. Planning across chart folds is outside this
project's core scope.

For every chart, keep pose validity, joint limits,
full-Jacobian regularity, and reduced-chart regularity as separate diagnostics.
The rectangular $6\times7$ Jacobian has no determinant. A feasible pixel is
not a validated edge: re-solve intermediate prescribed poses and check joint
continuity within limits. If collision checking is unavailable, state that
feasibility is kinematic only; a clearance proxy is not a full collision test.

\section*{Expected outcome}
Understand the iiwa model, fixed-angle IK classification, and how
feasibility varies over path progress and redundancy. Deliver branch-labelled
charts, a validated schedule where found, and a quantitative explanation of
the difference between constant and varying redundancy.

Submit reproducible code, model parameters and test data, the required plots,
a concise 5--6-page report, and a short simulation demonstration. Explain
both successful cases and failures; conclusions must follow the numerical
and geometric evidence rather than an expected outcome.

\section*{Resources}
The KUKA LBR iiwa 7 URDF is
\path{lectures_main/assets/models/iiwa7/iiwa7.urdf}; STL meshes are in that
directory. The accompanying \path{iiwa7_free_joints.urdf} is a mathematical
continuation model; apply physical limits separately. The course analytical
IK model will be supplied; reuse it rather than deriving a general 7R solver.

The course pose path is
\path{lectures_main/assets/data/path/kuka_iiwa7_rectangle_workspace_path_500_quaternion.csv}.
Its quaternion columns are \texttt{qw,qx,qy,qz}; normalize them and make their
signs consistent before interpolation. The reference table
\path{lectures_main/assets/data/null/iiwa7_IK_q3_030deg.csv}
checks fixed-$q_3$ solutions only; it does not replace evaluation at other angles.
Reuse Lecture 8 simulations and Exercise 4 data for connection validation
and schedule selection.

Use the lecture website to verify D--H parameters, visualize IKs and
singularities, and simulate paths. All listed paths are relative to the
repository root. Start the website following \path{lectures_main/README.md}.

\section*{Workload and collaboration}
Budget approximately 60 hours per student (120 person-hours per pair)
for this 2-ECTS project, following EPFL's 30-hours-per-credit convention.\footnote{\href{https://www.epfl.ch/education/teaching/teaching-guide-2/getting-started/design-a-course_1/}{EPFL: Design a Course, workload guidance}.}
Deduct any lectures or exercises included in those credits. Allocate roughly
20 team-hours to model validation, 50 to the core work, 30 to experiments,
and 20 to reporting. Reuse the specified IK and simulations. Share validation
and interpretation even if modelling and implementation are divided between
students. Hardware execution is outside scope.

\end{document}
